\documentclass[10pt,letterpaper]{article}
\usepackage{cornell-notes}

\title{Annex B.02.08: 02-08-correspondence-methods}
\author{}
\date{}
\setDocTitle{Annex B.02.08: 02-08-correspondence-methods}
\setDocSubtitle{ISO/IEC/IEEE 42010:2022}
\setDocOwner{ISO 42010 Cornell Notes}
\setDocDate{\today}
\setCornellCollection{ISO/IEC/IEEE 42010:2022 Cornell Notes}
\setCornellUnitType{Annex}
\setCornellUnitNumber{B.02.08}
\setCornellUnitTitle{02-08-correspondence-methods}
\hypersetup{pdftitle={Annex B.02.08: 02-08-correspondence-methods},pdfsubject={Cornell notes},pdfauthor={}}

\begin{document}
\maketitle
\begin{CornellOverviewBox}{Overview}
{\Large\bfseries\textcolor{CNavy}{Annex B.2.8: Correspondence methods}}\\[3pt]
{\small\textbf{Classification:} Informative annex}\\
{\small\textbf{Learning objective:} Use the informative guidance on correspondence methods to reinforce the normative and conceptual material without treating the guidance itself as mandatory.}
\end{CornellOverviewBox}
\vspace{3pt}

\begin{CornellNotesTable}
\CornellNoteRow{Purpose / key concept}{A listing of any correspondence methods defined by this viewpoint or its model kinds (per 8.1, 8.2 and 6.9.3).}
\CornellNoteRow{Core details}{\begin{itemize}[leftmargin=*,nosep,topsep=1pt]
\item These methods can be applied across view components, across views within an AD or across ADs.
\end{itemize}}
\CornellNoteRow{Related sections}{8.1, 8.2, 6.9.3}
\CornellNoteRow{Active recall}{\begin{itemize}[leftmargin=*,nosep,topsep=1pt]
\item What is the central role of Correspondence methods?
\item How does Correspondence methods connect to adjacent architecture-description concepts?
\end{itemize}}
\end{CornellNotesTable}


\begin{CornellSummaryBox}{Summary}
A listing of any correspondence methods defined by this viewpoint or its model kinds (per 8.1, 8.2 and 6.9.3). This material is informative guidance and does not itself create a conformance requirement.
\end{CornellSummaryBox}

\end{document}
